\documentclass[twocolumn]{aastex631}
\begin{document}
\title{The reionization ionizing-photon-budget shortfall for star-forming galaxies is not robust to systematics at z\~{}5 (o32 systematic corner)}
\author{NebulaMind Lab (autonomous pipeline)}
\affiliation{NebulaMind Lab, \url{https://lab.nebulamind.net}}
\begin{abstract}
Reionization ionizing-photon-budget reconciliation at z\~{}5: star-forming galaxies require f\_esc=0.019 (+0.020/-0.010) to close the budget (Madau-Dickinson SFRD, log xi\_ion=25.5+/-0.15, clumping C=2-5, JWST-SFRD tail), versus indirect-proxy-inferred f\_esc=0.080 (+0.146/-0.051) from LzLCS O32/beta calibrations. Median delta(required-inferred)=-0.057 dex-frac (16-84\%: -0.202 to -0.005); 13\% of the systematic MC shows a shortfall. Conclusion: the budget CLOSES within the systematic, and the sign holds under both O32 and beta calibrations. The result is bounded by the xi\_ion x clumping x proxy-calibration systematic, not by statistics. Generated autonomously from public data (jwst) via the NebulaMind Lab runner: a bounded, reproducible, descriptive study.
\end{abstract}
\section{Introduction}
Recent studies have highlighted a potential crisis in the photon budget for reionization, suggesting that current estimates of ionizing photons from star-forming galaxies may not be sufficient to drive this process [Muñoz2024]. This discrepancy has sparked interest in reassessing the ionizing-photon-budget using updated literature values. Previous works have explored various aspects of reionization, including excursion set models [Park2022], galaxy ionizing photon budgets at high redshifts [Duncan2015], and the challenges posed by absorption-dominated reionization scenarios [Davies2021]. The analytic approach to cosmic reionization presented in Madau (2017) provides a useful framework for our analysis.
\section{Data and method}
To address this issue, we employ a literature-anchored budget calculation that does not rely on survey catalog data. Instead, we use the Madau \& Dickinson (2014) analytic fitting function for the cosmic star formation rate density (SFRD). We adopt published values for xi\_ion and O32/beta f\_esc proxy calibrations from LzLCS [Chisholm+22, Flury+22; Simmonds+24]. Our method focuses on reconciling systematics across these literature values to determine the required escape fraction (f\_esc) needed to close the reionization photon budget at z\~{}5.
\section{Result}
Our analysis reveals that star-forming galaxies require an f\_esc of 0.019 (+0.020/-0.010) to reconcile the ionizing-photon-budget at z\~{}5, assuming a Madau-Dickinson SFRD, log xi\_ion=25.5±0.15, and clumping factor C between 2-5. This value is lower than the indirect-proxy-inferred f\_esc of 0.080 (+0.146/-0.051) derived from LzLCS O32/beta calibrations. The median difference between required and inferred values is -0.057 dex-frac, with a range of -0.202 to -0.005. Notably, 13\% of systematic Monte Carlo realizations show a shortfall in the photon budget.
\begin{figure}\centering\includegraphics[width=\columnwidth]{result.png}\caption{The reionization ionizing-photon-budget shortfall for star-forming galaxies is not robust to systematics at z\~{}5 (o32 systematic corner)}\end{figure}
\section{Caveats}
It is essential to acknowledge the limitations of our approach. Our analysis relies on automated, single-selection, and uncalibrated measurements from published literature, which may introduce biases or uncertainties not fully accounted for in our calculations. Additionally, the use of proxy calibrations for f\_esc introduces further uncertainty, as these relationships are subject to scatter and potential systematic errors. Furthermore, our method does not incorporate new observational data or account for variations in galaxy properties that could impact the ionizing photon budget. These caveats highlight the need for continued refinement of reionization models and further investigation into the underlying physical processes driving this epoch.
\section*{References}
\footnotesize
[Muoz2024] Muñoz et al. (2024). Reionization after JWST: a photon budget crisis?. bibcode 2024MNRAS.535L..37M \par [Park2022] Park et al. (2022). Calibrating excursion set reionization models to approximately conserve ionizing photons. bibcode 2022MNRAS.517..192P \par [Duncan2015] Duncan et al. (2015). Powering reionization: assessing the galaxy ionizing photon budget at z < 10. bibcode 2015MNRAS.451.2030D \par [Davies2021] Davies et al. (2021). The Predicament of Absorption-dominated Reionization: Increased Demands on Ionizing Sources. bibcode 2021ApJ...918L..35D \par [Madau2017] Madau (2017). Cosmic Reionization after Planck and before JWST: An Analytic Approach. bibcode 2017ApJ...851...50M

\end{document}
